\documentclass[11pt,a4paper]{article}
% =====================================================================
%  Shared preamble for the Part V lecture notes (lectures 19-22).
%
%  These four lectures sit outside Geron, so the notes ARE the primary
%  source and are examined on the same terms as the textbook parts.
%  Everything here is chosen to make that readable on paper: one column,
%  generous measure, and the same six-colour palette as the slides so a
%  student moving between the two is not re-learning what red means.
%
%  Build with pdflatex (twice, for the toc and the refs):
%      cd notes && latexmk -pdf lecture-19.tex
%  or  make -C notes
% =====================================================================
\usepackage[T1]{fontenc}
\usepackage[utf8]{inputenc}
\usepackage{newtxtext,newtxmath}      % Times-like: prints well, reads well
\usepackage[scaled=0.86]{inconsolata} % code, at a size that sits beside Times
\usepackage{microtype}
% newtx defines \Bbbk, \openbox, \proof and \endproof; amssymb and amsthm
% then redefine them and abort. Freeing the four names before those two
% packages load is the standard fix, and it costs nothing here: we use
% amsthm only for \newtheorem.
\let\Bbbk\relax
\usepackage{amsmath,amssymb}
\let\openbox\relax
\let\proof\relax
\let\endproof\relax
\usepackage{amsthm}
\usepackage{booktabs}
\usepackage{enumitem}
\usepackage{xcolor}
\usepackage{tcolorbox}
\usepackage{titlesec}
\usepackage{graphicx}
\usepackage[hidelinks]{hyperref}
\tcbuselibrary{skins,breakable}

% ---- the course palette, from assets/css/custom.css -------------------
\definecolor{aimlprimary}{HTML}{0B3D62}   % deep blue  - structure
\definecolor{aimlaccent} {HTML}{C0392B}   % brick red  - failures
\definecolor{aimlsuccess}{HTML}{14663A}   % green      - repairs
\definecolor{aimlmath}   {HTML}{6C3483}   % purple     - the derivation
\definecolor{aimlmuted}  {HTML}{4B5563}
\definecolor{aimlrule}   {HTML}{B0BCC7}
\definecolor{aimlsoft}   {HTML}{F4F7F9}

\usepackage[a4paper,margin=32mm,top=30mm,bottom=30mm]{geometry}
\setlength{\parskip}{0.45\baselineskip}
\setlength{\parindent}{0pt}
\linespread{1.06}

% ---- headings --------------------------------------------------------
\titleformat{\section}{\Large\bfseries\color{aimlprimary}}{\thesection}{0.7em}{}
\titleformat{\subsection}{\large\bfseries\color{aimlprimary}}{\thesubsection}{0.7em}{}
\titleformat{\subsubsection}{\bfseries\color{aimlmuted}}{}{0em}{}
\titlespacing*{\section}{0pt}{1.6\baselineskip}{0.5\baselineskip}
\titlespacing*{\subsection}{0pt}{1.2\baselineskip}{0.35\baselineskip}

% ---- boxes -----------------------------------------------------------
% One box per role, and the role is the colour. A student who has seen the
% slides already knows what each one means before reading a word of it.
\newtcolorbox{keybox}[1][]{
  breakable, enhanced, colback=aimlsoft, colframe=aimlprimary,
  boxrule=0.4pt, left=8pt, right=8pt, top=6pt, bottom=6pt,
  fonttitle=\bfseries\small, coltitle=white, colbacktitle=aimlprimary,
  attach boxed title to top left={yshift=-2mm, xshift=4mm},
  boxed title style={boxrule=0pt, arc=1pt}, #1}

\newtcolorbox{failbox}[1][]{
  breakable, enhanced, colback=white, colframe=aimlaccent,
  boxrule=0.9pt, left=8pt, right=8pt, top=6pt, bottom=6pt,
  fonttitle=\bfseries\small, coltitle=white, colbacktitle=aimlaccent,
  attach boxed title to top left={yshift=-2mm, xshift=4mm},
  boxed title style={boxrule=0pt, arc=1pt}, #1}

\newtcolorbox{derivbox}[1][]{
  breakable, enhanced, colback=white, colframe=aimlmath,
  boxrule=0.6pt, left=8pt, right=8pt, top=6pt, bottom=6pt,
  fonttitle=\bfseries\small, coltitle=white, colbacktitle=aimlmath,
  attach boxed title to top left={yshift=-2mm, xshift=4mm},
  boxed title style={boxrule=0pt, arc=1pt}, #1}

% An exercise. Deliberately the same shape as the notebook's `try` field:
% one change, and what should happen to the output.
\newtcolorbox{trybox}[1][]{
  breakable, enhanced, colback=white, colframe=aimlsuccess,
  boxrule=0.5pt, left=8pt, right=8pt, top=6pt, bottom=6pt,
  fonttitle=\bfseries\small, coltitle=white, colbacktitle=aimlsuccess,
  attach boxed title to top left={yshift=-2mm, xshift=4mm},
  boxed title style={boxrule=0pt, arc=1pt}, #1}

% ---- small semantics -------------------------------------------------
\newcommand{\scope}[1]{\textcolor{aimlmuted}{\small #1}}
\newcommand{\term}[1]{\textbf{#1}}
\newcommand{\code}[1]{\texttt{\small #1}}
\newcommand{\lecture}[1]{Lecture~#1}

\newtheorem{definition}{Definition}[section]
\newtheorem{proposition}[definition]{Proposition}

% ---- title block -----------------------------------------------------
% \notestitle{19}{Information retrieval: the lexical foundation}{SciFact (BEIR)}
\newcommand{\notestitle}[3]{%
  \thispagestyle{empty}
  \begin{flushleft}
    {\color{aimlmuted}\small\sffamily
     APPLICATIONS OF MACHINE LEARNING \quad$\cdot$\quad
     BSc Mathematics of Artificial Intelligence}\\[2mm]
    {\color{aimlrule}\rule{\textwidth}{0.8pt}}\\[4mm]
    {\color{aimlmuted}\large Lecture #1 \quad$\cdot$\quad Part V \quad$\cdot$\quad
     Extended notes}\\[3mm]
    {\Huge\bfseries\color{aimlprimary}\baselineskip=1.15\baselineskip #2\par}\vspace{4mm}
    {\color{aimlmuted}\large Dataset: #3}\\[3mm]
    {\color{aimlrule}\rule{\textwidth}{0.8pt}}
  \end{flushleft}
  \vspace{2mm}
  \begin{keybox}[title={Why these notes exist}]
    Lectures 19--22 sit outside G\'eron, the course's primary text. For those
    four lectures \emph{these notes are the primary source}, and they are
    examined on exactly the same terms as the textbook parts. Everything in
    the slide deck for this lecture is here, with the arguments written out at
    length rather than compressed onto a slide; the numbers are the ones the
    accompanying notebook prints, and every one of them is a measurement on
    the corpus named above rather than a figure quoted from a paper.
  \end{keybox}
  \vspace{1mm}
  {\small\tableofcontents}
  \vspace{2mm}
  \noindent{\color{aimlrule}\rule{\textwidth}{0.4pt}}
  \vspace{1mm}
}


\begin{document}
\notestitlefor{10}{II}{PyTorch}{Fashion-MNIST}{Chapter 10}

\section{Where we are}

Lecture 9 found four things \code{MLPClassifier} will not let you do: change
what it optimises, see a gradient, stop mid-epoch, or move the computation to
another device. All four are one wall --- \code{fit()} contains a loop over
epochs and batches that you did not write, cannot see and cannot enter.

The repair is not a better library. It is \emph{owning the loop}. Today: what
that loop has to compute, why it can be computed at all, and how to write it.

\section{Derivation: backpropagation as reverse-mode autodiff}

You called \code{clf.fit()}. It computed 266{,}610 partial derivatives, on
every batch, twenty times over. \emph{How?} The cost of that computation is
what decides whether training a network is possible at all --- and the reason
it is possible is a property of the \emph{direction} in which the chain rule is
applied.

\begin{center}
\small
\begin{tabular}{llp{0.28\textwidth}}
\toprule
\textbf{Method} & \textbf{Cost for $P$ parameters} & \textbf{Exact?} \\
\midrule
By hand & your afternoon, per architecture & yes, if you are careful \\
Finite differences & $P+1$ evaluations & no --- truncation and cancellation \\
Symbolic & expression blow-up & yes \\
Automatic differentiation & see below & to machine precision \\
\bottomrule
\end{tabular}
\end{center}

Automatic differentiation is neither symbolic nor numerical. It differentiates
the \emph{program}, operation by operation, using the chain rule and nothing
else.

\subsection{Two bracketings, two algorithms}

A network is a composition, so writing
$L = f_3(f_2(f_1(\boldsymbol\theta)))$ gives
\[ \frac{\partial L}{\partial\boldsymbol\theta} = J_3 J_2 J_1, \]
with $J_i$ the Jacobian of $f_i$ evaluated at the point the forward pass
reached. Every layer contributes one factor --- and matrix multiplication is
associative:
\[ \underbrace{(J_3 J_2) J_1}_{\text{reverse mode}}
   \;=\;
   \underbrace{J_3 (J_2 J_1)}_{\text{forward mode}}. \]
Same number. Wildly different cost.

\term{Forward mode} picks a direction $\mathbf{v}$ in input space and carries a
tangent $\dot u = \partial u/\partial v$ alongside every intermediate value.
One sweep gives the derivative of every output with respect to \emph{one} input
direction, and nothing needs to be stored, because the tangent travels with the
value.

\term{Reverse mode} picks a direction in output space --- for a scalar loss
there is only one, and it is 1 --- and carries an adjoint
$\bar u = \partial L/\partial u$ backwards. One sweep gives the derivative of
\emph{one} output with respect to every input, and the forward values must be
kept, because the Jacobians are evaluated at them.

\begin{center}
\small
\begin{tabular}{lll}
\toprule
\textbf{Mode} & \textbf{Sweeps needed} & \textbf{Cheap when} \\
\midrule
Forward & $n$ --- one per input & few inputs, many outputs \\
Reverse & $m$ --- one per output & many inputs, few outputs \\
\bottomrule
\end{tabular}
\end{center}

Each sweep costs a small constant times one forward evaluation --- around 2 to
4, and it does not depend on $n$ or $m$. That constant is why the backward pass
costs about what the forward pass costs, no matter how many parameters there
are.

\subsection{A worked example small enough to check}

\[ L = w^2, \qquad w = xy + \sin x, \qquad \text{at } x = 2,\ y = 3. \]
By hand, $\partial L/\partial x = 2w(y + \cos x)$ and
$\partial L/\partial y = 2wx$.

\begin{center}
\small
\begin{tabular}{llr@{\hskip 2.2em}llr}
\toprule
\textbf{Node} & \textbf{Expression} & \textbf{Value} &
\textbf{Adjoint} & \textbf{Computed as} & \textbf{Value} \\
\midrule
$x$ & input      & 2.0    & $\partial L/\partial L$ & the seed & 1 \\
$y$ & input      & 3.0    & $\partial L/\partial w$ & $2w$ & 13.8186 \\
$u$ & $x\cdot y$ & 6.0000 & $\partial L/\partial u$ & $(\partial L/\partial w)\cdot 1$ & 13.8186 \\
$s$ & $\sin x$   & 0.9093 & $\partial L/\partial s$ & $(\partial L/\partial w)\cdot 1$ & 13.8186 \\
$w$ & $u + s$    & 6.9093 & $\partial L/\partial y$ & $(\partial L/\partial u)\,x$ & 27.6372 \\
$L$ & $w^2$      & 47.7384 & $\partial L/\partial x$ & $(\partial L/\partial u)\,y + (\partial L/\partial s)\cos x$ & 35.7052 \\
\bottomrule
\end{tabular}
\end{center}

Six nodes evaluated once each in topological order, then six lines backwards
--- one pass, both derivatives. Forward mode would need two sweeps here, one
seeded at $x$ and one at $y$.

\begin{keybox}[title={The only bookkeeping in the whole method}]
$x$ is used twice, once by the product and once by the sine, so
\[ \frac{\partial L}{\partial x}
   = \underbrace{41.4557}_{\text{via }\times}
     \underbrace{-\,5.7505}_{\text{via }\sin}
   = 35.7052. \]
A node's adjoint is the \emph{sum} of the contributions along its outgoing
edges. Miss one edge and the gradient is silently wrong, not absent.

Reverse-mode autodiff is a topological sort, a table of local derivatives, and
this sum. There is no other idea in it.
\end{keybox}

Checked against the library, \code{torch} returns
$(35.705220,\ 27.637190)$ --- agreeing with the hand derivatives to machine
precision. The same habit as Lecture 2's orthogonality test.

\subsection{The argument that settles it}

Training differentiates a function from 266{,}610 parameters to \emph{one}
scalar loss. Forward mode needs one sweep per input: 266{,}610 of them. Reverse
mode needs one per output: 1.

\begin{keybox}[title={Reverse mode is not a good choice}]
It is the \emph{only viable} one, by a factor equal to the parameter count.
\end{keybox}

And the loss being one number is not a convention. Gradient descent needs a
direction of steepest descent, which is only defined for a scalar-valued
function; a vector-valued objective has a Jacobian, not a gradient, and no
canonical direction to step in. So the shape of the training problem --- many
parameters, one number --- is fixed by the optimisation, and that shape is
precisely the one reverse mode is cheap for. The two facts are not independent.

\subsection{Measure it rather than trusting it}

Small networks, where forward mode can actually be run to completion. Same
loss, same point, both modes --- and the two columns agree to
$2.235\times10^{-8}$, so this is a measurement of cost, not of quality.

\begin{center}
\small
\begin{tabular}{rrrr}
\toprule
\textbf{Parameters} & \textbf{Reverse, whole gradient} &
\textbf{Forward, whole gradient} & \textbf{Ratio} \\
\midrule
 51 & 0.24 ms &   9.3 ms &  39$\times$ \\
 99 & 0.22 ms &  16.8 ms &  76$\times$ \\
195 & 0.26 ms &  35.5 ms & 137$\times$ \\
387 & 0.21 ms &  73.1 ms & 341$\times$ \\
771 & 0.27 ms & 151.7 ms & 572$\times$ \\
\bottomrule
\end{tabular}
\end{center}

Reverse mode is flat in $P$; forward mode is linear. At the size we actually
train, reverse mode costs 1.33 ms for the whole gradient and forward mode 0.78
ms for \emph{one direction} --- so all 266{,}610 directions would take about
3.5 minutes. That last figure is projected, not run, and saying so is the
honest form of the claim.

Lecture 9's run took 20 epochs of 430 batches: 8{,}600 gradients, in 125
seconds. The same 8{,}600 in forward mode, at the measured per-pass cost, would
be \textbf{21 days}. The associativity of matrix multiplication is the only
reason this course exists.

\begin{keybox}[title={So what is backpropagation?}]
Reverse-mode automatic differentiation, applied to a neural network, with the
seed set to 1 at the loss. It is not specific to neural networks, it is not an
approximation, and it is not gradient descent --- it supplies the gradient that
descent then uses. The 1986 paper's contribution was noticing this was cheap
enough to be practical, not inventing the chain rule.
\end{keybox}

\begin{failbox}[title={What reverse mode costs}]
Memory. Every Jacobian is evaluated at the point the forward pass reached, so
every intermediate activation must be kept until the backward pass consumes it.
Forward mode stores nothing; reverse mode stores the whole forward trace, and
the trace grows with the batch size and with the depth.

Lecture 12 computes the RAM of a convolutional network and finds that the
activations, not the parameters, are what exhausts the device. This is why.
\end{failbox}

\section{What the wrapper costs, measured}

The accuracy was not the bug. 88.02\% against a 10.0\% baseline is a real
result: the split was clean, the validation set separate, the test set touched
once. The bug is that you cannot do anything else.

\begin{center}
\small
\begin{tabular}{lrr}
\toprule
\textbf{Same architecture, optimiser and 20 epochs} & \textbf{Wall clock} &
\textbf{Test accuracy} \\
\midrule
Scikit-Learn, CPU & 125 s & 88.02\% \\
PyTorch, CPU      &  14 s & 88.84\% \\
PyTorch, MPS      &  14 s & 89.24\% \\
\bottomrule
\end{tabular}
\end{center}

9.1$\times$ faster for the same model --- and look at rows two and three. The
accelerator bought 1.00$\times$, that is, nothing. \emph{All of the speed-up is
the framework and none of it is the device.}

The 9.1$\times$ came from how the arithmetic is dispatched. Scikit-Learn's MLP
is written for clarity and generality, in NumPy, with a Python-level loop over
batches; PyTorch fuses the same operations into fewer, larger kernels and keeps
the tensors in one layout throughout. Same architecture, same optimiser, same
epochs, same arithmetic --- and 125 seconds against 14, \emph{both on the CPU}.

\subsection{So when does the accelerator pay?}

A forward pass is $\phi(\mathbf{X}\mathbf{W} + \mathbf{b})$ repeated, and a GPU
is a machine for large matrix products and very little else. Ours are not
large: 266{,}610 parameters and a batch of 128 is a matrix product that
finishes before the cost of dispatching it has been paid.

\begin{center}
\small
\begin{tabular}{lrrrr}
\toprule
\textbf{Hidden layers} & \textbf{Parameters} & \textbf{CPU} & \textbf{MPS} &
\textbf{Ratio} \\
\midrule
100        &    79{,}510 & 0.13 s & 0.45 s & 0.28$\times$ \\
300, 100   &   266{,}610 & 0.21 s & 0.25 s & 0.84$\times$ \\
1000, 1000 & 1{,}796{,}010 & 0.62 s & 0.28 s & 2.22$\times$ \\
2000, 2000 & 5{,}592{,}010 & 1.03 s & 0.43 s & 2.40$\times$ \\
\bottomrule
\end{tabular}
\end{center}

At our size the ratio is \emph{below} 1: the accelerator is slower. It starts
winning at about a million parameters, once each batch carries enough
arithmetic to cover the dispatch. Measured on an Apple MPS backend; a Colab T4
has a different crossover, in the same place in the argument.

\begin{keybox}[title={The honest version of ``you need a GPU''}]
You need one when the network is large enough for the device to matter, and not
before. Nothing in this course is --- every notebook is sized for Colab's free
CPU tier, deliberately, so that every number on these pages is one you can
reproduce.

Saying ``you need a GPU'' without measuring where the crossover falls would be
exactly the sort of unmeasured claim this course exists to stop.
\end{keybox}

\section{Tensors, autograd, and the loop}

A tensor is an array that also knows two things a NumPy array does not: where
it lives --- an operation between tensors on different devices is an error, and
a loud one --- and whether to record.

\code{requires\_grad} does not compute a derivative. It says \emph{record what
happens to me}. Every operation on a recording tensor creates a node holding
the operation and its inputs; the graph is built during the forward pass, by
running it, with no separate compilation step; and it is rebuilt from scratch
on every batch, which is why a Python \code{if} inside \code{forward} just
works. That is the forward sweep from \S2, with the graph written down as it
goes so the reverse sweep has something to walk --- and you can inspect it:
\code{r.grad\_fn} and \code{r.grad\_fn.next\_functions} are the nodes of the
worked example, by name.

Evaluation needs no derivatives, so \code{torch.no\_grad()} stops the recording
and saves memory proportional to the batch. Forgetting it is slow and
memory-hungry, but it is not a correctness bug. The two that follow are.

\begin{center}
\small
\begin{verbatim}
for epoch in range(EPOCHS):
    model.train()
    for xb, yb in train_loader:
        opt.zero_grad()          # 1  clears p.grad, which backward() adds to
        out  = model(xb)         # 2  the forward sweep; the graph is recorded
        loss = lossf(out, yb)    # 3  reduces the batch to one scalar
        loss.backward()          # 4  the reverse sweep; fills p.grad
        opt.step()               # 5  updates the parameters from p.grad
\end{verbatim}
\end{center}

Five lines. There is no hidden machinery: this is the whole of what
\code{fit()} was doing --- and every arrow is a line you wrote, so every arrow
is a place you can print, clip, log, branch or stop.

Rebuilt in PyTorch the model scores 89.24\% against Scikit-Learn's 88.02\% in
14 seconds against 125. The accuracies agree, as they must: same architecture,
same optimiser, same learning rate, same epochs, same data. \emph{The rebuild
is not an improvement --- it is the same model.} What we bought is the loop.
Overlaying the two learning curves is the first check after any rewrite; if
they had disagreed, one of the two implementations would be wrong.

\section{Two bugs that never raise}

\subsection{The missing \code{zero\_grad()}}

\code{backward()} \emph{accumulates}: \code{p.grad} is added to, not
overwritten. After one call the gradient norm is 1.6183, after two 3.2366,
after three 4.8549, and after \code{zero\_grad()} it is 0.0000.

That behaviour is deliberate and useful --- it is how you accumulate a large
effective batch across several small ones, and how you sum gradients from two
loss terms. It is only a bug when you did not ask for it.

\begin{failbox}[title={What it does to training}]
At step $t$ the optimiser sees the sum of every gradient computed so far, not
the current one. The effective step size grows without bound, and Adam's
normalisation hides the divergence for a while. No exception, no warning, no
\code{NaN}.

\begin{center}
\small
\begin{tabular}{lrr}
\toprule
& \textbf{Final training loss} & \textbf{Validation accuracy} \\
\midrule
With \code{opt.zero\_grad()} & 0.2574 & 87.94\% \\
Without it                   & 2.2933 & 11.20\% \\
\bottomrule
\end{tabular}
\end{center}
76.7 points, on identical seeds and identical everything else.
\end{failbox}

To catch it: read the loop for the five lines in order --- a checklist item,
not an insight --- and print the gradient norm for the first few steps, since
if it climbs monotonically it is accumulating. Watch \emph{where} it sits, too:
\code{zero\_grad()} in the epoch loop instead of the batch loop is the same
bug, indented two spaces differently, and it survives review.

\begin{keybox}[title={The general habit}]
When a framework's default is ``accumulate'', the code that does \emph{not}
accumulate is the code with an extra line in it. Ask what happens if that line
is deleted --- for every line in the loop.
\end{keybox}

\subsection{The missing \code{model.eval()}}

Some layers behave differently in the two phases. Dropout zeroes a random
fraction of units in \code{train()} and passes everything through in
\code{eval()}; batch normalisation uses this batch's statistics in one and the
running averages in the other. \emph{The module defaults to training mode}, so
building a model and never calling \code{eval()} means every evaluation you
make is of a randomly crippled network.

\begin{center}
\small
\begin{tabular}{lr}
\toprule
\textbf{Evaluation mode} & \textbf{Test accuracy} \\
\midrule
\code{model.eval()} & 86.77\% \\
\code{model.train()}, mean of 10 calls & 86.26\% \\
\quad lowest of those 10 & 86.06\% \\
\quad highest of those 10 & 86.68\% \\
\bottomrule
\end{tabular}
\end{center}

\begin{keybox}[title={The tell is the spread, not the mean}]
0.62 accuracy points across ten \emph{identical} calls. A deterministic
function of fixed weights and fixed data returns the same number every time, so
if your metric wobbles between calls, ask which layer is still in training
mode.

Run any evaluation twice. It costs one line, and it catches this bug, the
seeding bugs, and any accidental shuffling of the labels.
\end{keybox}

Both bugs are one line, neither raises, neither warns, and both produce a
plausible number --- costing 76.7 points and 0.51 points respectively. One is
caught by reading the loop, the other by running it twice, and both survive a
review that is looking for exceptions rather than for behaviour.

\begin{failbox}[title={``Running it twice'' needs one caveat}]
A fixed seed does not make a GPU run bit-reproducible. cuDNN picks convolution
algorithms by benchmarking them, and reductions that accumulate with atomics
add their terms in whatever order the blocks finish; neither is seeded, so two
runs of identical code differ in about the third decimal.
\code{torch.use\_deterministic\_algorithms(True)} buys determinism back at a
cost in speed.

So calibrate what counts as evidence: without those flags a difference in the
third decimal means nothing. The two bugs above cost 76.7 and 0.51 points ---
both far outside it.
\end{failbox}

\section{Feeding the loop, and one more silent error}

\code{Dataset} answers two questions --- how many, and give me number $i$ ---
and \code{DataLoader} shuffles, batches, and can load in parallel while the
device is busy. Indexing a big tensor by hand works while the data fits in
memory; it stops working in Lecture 12, and the interface does not change.

\begin{failbox}[title={The last batch is short}]
10{,}000 test images in batches of 384 gives 27 batches, the last containing
16 --- so a plain mean of per-batch accuracies weights those 16 exactly as
heavily as a full batch.

\begin{center}
\small
\begin{tabular}{lr}
\toprule
\textbf{How the accuracy was aggregated} & \textbf{Value} \\
\midrule
Counted over the whole set & 89.240\% \\
Mean of the per-batch accuracies & 89.622\% \\
\bottomrule
\end{tabular}
\end{center}

0.382 points --- small \emph{here}, because the last batch is 16 of 384 and the
model is no better or worse on it. Change either and it grows: with batch size
3{,}000 on 10{,}000 images the short batch carries a quarter of the weight
instead of a tenth.

\code{drop\_last=False} is the default and that is the right default. The bug is
in how the metric is aggregated, not in the loader.
\end{failbox}

Weight by the batch size, or count over the set --- never take a plain mean of
per-batch metrics. Not because the error is always large, but because
\emph{you cannot tell from the number whether it was}. That is the identical
argument, in the identical shape, as Lecture 2's scale-before-split leak, where
the damage was nothing measurable and the rule stayed procedural for the same
reason. And it is worse for losses than for accuracy: a per-batch mean of a
per-batch mean is two averagings deep.

Subclassing \code{nn.Module} is what you reach for when \code{Sequential} runs
out --- two inputs or two outputs, a skip connection, anything conditional on
the data. All three are ordinary Python inside \code{forward}, because the
graph is rebuilt by running the function on every batch. One rule:
\code{super().\_\_init\_\_()} before anything else, or parameter registration
silently does not happen.

\section{The notebook}

\begin{trybox}[title={What to change}]
The notebook prints which device it found in its first cell. Leave it on
\code{cpu}: everything here is sized for the free tier, and the crossover
measured in \S3 shows the accelerator running \emph{this} network slower.
Change it only to reproduce that measurement --- and then you will have
measured your own crossover rather than accepted mine.
\end{trybox}

\section{Exercises}

\begin{enumerate}[leftmargin=1.6em,itemsep=4pt]
  \item The chain rule for a composition is a product of Jacobians. Explain why
        bracketing it from the output end costs one sweep per output, and state
        which end is right for training a network. \hfill \scope{[5]}
  \item In the worked example $L = (xy + \sin x)^2$, the adjoint of $x$ is a
        sum of two terms. Say why, and what goes wrong if one is omitted.
        \hfill \scope{[4]}
  \item A run reports 9.1$\times$ speed-up moving from Scikit-Learn to PyTorch,
        and 1.00$\times$ moving from CPU to an accelerator. Explain both
        numbers. \hfill \scope{[4]}
  \item A training loop omits \code{opt.zero\_grad()}. State what the optimiser
        sees at step $t$, why no exception is raised, and one diagnostic that
        would reveal it in the first few steps. \hfill \scope{[5]}
  \item An evaluation returns 86.26\%, 86.68\% and 86.06\% on three identical
        calls with fixed weights and fixed data. Name the cause and the
        one-line fix. \hfill \scope{[4]}
\end{enumerate}

\section*{Summary}
\addcontentsline{toc}{section}{Summary}

The chain rule is a product of Jacobians and the product may be bracketed from
either end. Forward mode costs one sweep per input; reverse mode costs one
sweep per output; training has 266{,}610 inputs and one scalar output, so
reverse mode is the only viable choice --- by a factor equal to the parameter
count. Measured, forward mode is linear in $P$ and reverse mode flat, and
Lecture 9's 8{,}600 gradients would take 21 days instead of 125 seconds. The
price is memory: the forward trace must be kept, which is the bill Lecture 12
pays.

Moving to PyTorch bought 9.1$\times$ and the accelerator bought nothing,
because at 266{,}610 parameters the matrix product finishes before its dispatch
is paid for --- the crossover is near a million. The rebuilt model is the
\emph{same} model, and what was bought is the five-line loop.

Then two bugs that raise nothing, warn nothing, and return a plausible number.
A missing \code{zero\_grad()} costs 76.7 accuracy points because
\code{backward()} accumulates by design; a missing \code{model.eval()} costs
0.51, and its tell is a 0.62-point spread across ten identical calls. One is
caught by reading the loop, the other by running it twice.

\end{document}
